\section{Policy extensions}\label{si:policy}

\paragraph{Half-Earth and reservation.} The Half-Earth proposal \citep{Wilson2016} caps human-usable biocapacity at half of the total. In the two-land setting this is a reservation $E\le\sigma B$ with $\sigma$ the human-available share, enforced either as a cap on the fast book or on the carrying capacity; with $\sigma_f,\sigma_c<1$ in Eq.~(4) the shortfall is computed against the reserved flows, so conversion begins earlier for a given demand and the ceiling of Proposition~2 is lowered to $\sigma$-weighted values. On the nominal path a reservation that keeps $E\le\sigma B$ succeeds with no liquidation. This is a nominal-path consequence until it is certified against a declared disturbance class (persistent productivity, $b_f$ or $\rho_c$ shocks) under a preregistered retention rule (a governance module is retained only if it improves a declared protection-and-supply score), with the model's declared fit defects converted into an erosion margin on the closed loop's contraction rate; until then it is a property of the nominal calibration, not a robustness claim.

\paragraph{Typed floors and freeze conversion.} Reserve $A_r\ge A_r^{\min}$, capital viability $A_c\ge\Acmin$, arable ceiling $A_f\le A_f^{\max}$ and $K\ge K_{\min}$ are separate constraints. Freezing conversion (capping $A_f$ at $A_f^{\max}$ or flooring $A_c$ at $\Acmin$) is the single operator that guarantees $A_c\ge\Acmin$; the conversion rate saturates as $A_c\to\Acmin$ by construction of Eq.~(5). A policy that observes only the aggregate $B$ cannot guarantee capital-land integrity: the failure is informational (aggregation relative to the typed floor), not dynamical.

\paragraph{Institutional signature.} In the mobilising/protective taxonomy of the delay companion \citep{Abaee2026a}, freeze-conversion and surplus-gated restoration are protective; the mobilising sign (demand responding to a perceived shortfall by converting more) is the destabilising one; deficit-gated restoration is not unambiguously protective because it can deepen the short-run shortfall (Section~\ref{si:gates}).

\paragraph{Weak and strong sustainability reading.} In the material-cycle reading of the ledger companion \citep{Abaee2026c}, weak sustainability is the idealised regime in which substitution and regeneration jointly redistribute matter so that it is used as it arises -- substitution admissible only through an identified physical pathway, never merely because a weighted aggregate permits it -- and strong sustainability is ``the regime in which that closure fails -- either because no identified physical pathway re-routes the extracted matter, or because depletion outruns the rate at which such a pathway could be deployed'', so that a separately binding floor on a critical stock binds. The composite $B$ is the weak-sustainability index: the linear, perfectly substitutable member of the aggregator family of the assessment companion \citep[Lemma~A of][]{Abaee2026b}; it can satisfy its aggregate floor while the typed floor on $A_c$ is violated by exactly the amount that the fast-land yield gains against the capital loss, and in that family no positive-elasticity aggregator is uniformly safe (its Theorem~S2(ii)). Under exact-tube semantics the intermediate passage during which $A_c$ falls is itself a violation even if the endpoint recovers.

\paragraph{Reporting rule.} A growth headline for total biocapacity carries no information about whether cropland growth is compensating for stagnation or offsetting changes in the other components. The informative assessment states, at every level at which an aggregate is used: component endpoint ratios, log contributions, attribution shares, data vintage and endpoint sensitivity.

\section{The one-stock comparator}\label{si:onestock}

Earlier versions of this work developed a single-stock model that the two-land model supersedes. It is retained here as a comparator so that quantities that circulate for it are not mistaken for two-land results. The comparator has three states ($A$, $P$, $D$) and one algebraic relation:
\begin{align}
G(A)&=\rho A(1-A/A_{\max}),\qquad B=bA+b_GG(A),\qquad E=eP,\\
\dot A&=G\big(A(t-\tau_g)\big)-\pos{E-\sigma bA}/b_G,\\
\dot P&=rP\big(1-P/K(t-\tau_p)\big),\qquad K=B/e,\\
\dot D&=\pos{E-B}-\eta D,\qquad b=(b_0+T_b)e^{-\alpha D},
\end{align}
with demand met first from the flow $bA$ and the shortfall liquidating the standing stock at rate $1/b_G$. Its properties, each of which was re-derived under the two-land structure and found not to carry over, are the following.
\begin{itemize}[leftmargin=1.5em]
\item \emph{MSY and fold.} $B(A)$ has an interior maximum (the MSY of Schaefer, 1954) only in the capital-dominated regime $b_G\rho>b$; a fixed liability above $B_{\max}$ produces a saddle-node, the catastrophic-shift structure of \citet{Scheffer2001}. At the comparator baseline ($b_G\rho=0.04<b=0.5$) the fold is absent. In the two-land model $Y_c$ is linear in $A_c$ and $B_{eq}(A_c)$ is monotone; there is no fold in any regime (Proposition~2).
\item \emph{The equilibrium family $P=B(A)/e$ and $\det J\equiv0$.} Under full harvest ($\sigma=1$) the comparator has a neutral continuum of equilibria on which the Jacobian is exactly singular. The two-land model has a one-parameter family in the \emph{land allocation} (indexed by $A_c$ on the no-conversion face, with $A_r/A_f=\eta_f/\mu$), not a free $P$, and no identically singular Jacobian.
\item \emph{The $+0.6$ eigenvalue and the two-gain Hopf classification.} The comparator's deficit-region linearisation has a positive real root ($a_{11}=G'(A^*)+b/b_G\approx0.608$--$0.625$ at its baseline) and a two-gain rule $\chi=q/(\rho-2q)$ selecting which delay destabilises ($\tau_g$-only for $\chi>1$, $\tau_p$-only for $\chi<1$, with critical delays $\tau_m\approx85.4$ and $\tau_p\approx231$\,yr). None of this is a two-land quantity: the two-land deficit-side spectrum is real and negative ($-0.054$, $-1.0\,\mathrm{yr}^{-1}$) and constant along the branch.
\item \emph{The regeneration cliff and the masking window.} The comparator shows a recover/collapse boundary at $\tau_g\approx18$--$20$\,yr and a productivity-illusion window of $\approx5.4$\,yr ($A:1.00\to0.86$ while $B:0.578\to0.647$; the mask vanishes for deficits $\gtrsim0.075$). In the two-land model the boundary is a demand threshold insensitive to $\tau_g$ (Table~\ref{tab:si_delay}) and the composition window is set by $\lambda_1$ and $\dconv$ (paper, Section~3.1). The recover fractions $0.54\to0.52\to0.50$ for $\tau_g=5\to25\to40$\,yr quoted in earlier versions were computed with an area-based intermediate formulation and are superseded by the identical-to-$10^{-3}$ result of Table~\ref{tab:si_delay}.
\end{itemize}
Figures~\ref{fig:si_s5}--\ref{fig:si_s12} reproduce four comparator figures for reference; the complete comparator figure set (S1--S15 of the deposited supplementary folder) remains in the repository.

\begin{figure}[htbp]
\centering
\includegraphics[width=0.8\textwidth]{S5_demo_overshoot_run.png}
\caption{One-stock comparator: representative overshoot run (stock $A$, population $P$, biocapacity $B$ and footprint $E$). Comparator only; not a two-land result.}
\label{fig:si_s5}
\end{figure}

\begin{figure}[htbp]
\centering
\includegraphics[width=0.8\textwidth]{S6_masking_window.png}
\caption{One-stock comparator: the productivity-illusion window ($A$ falls from $1.00$ to $0.86$ while $B$ rises from $0.578$ to $0.647$; window $\approx5.4$\,yr; the mask vanishes for deficits $\gtrsim0.075$). The two-land composition window of the paper (Fig.~1a) is a different object with a different mechanism.}
\label{fig:si_s6}
\end{figure}

\begin{figure}[htbp]
\centering
\includegraphics[width=0.8\textwidth]{S7_recovery_vs_collapse.png}
\caption{One-stock comparator: recovery versus collapse as a function of the regeneration delay (the $\approx18$--$20$\,yr cliff). Not present in the two-land model.}
\label{fig:si_s7}
\end{figure}

\begin{figure}[htbp]
\centering
\includegraphics[width=0.8\textwidth]{S12_basin_baseline.png}
\caption{One-stock comparator: basin of recovery at the comparator baseline. In the two-land model the corresponding object is the continuous drawdown-depth surface of Fig.~1c and Table~\ref{tab:si_drawdown}.}
\label{fig:si_s12}
\end{figure}

\section*{Supplementary references}
